%% ==========================================================================
%%  Approach 8 --- the global route: which optimum is good?
%%
%%  Source: derived 2026-08-06.  Scripts: update_20/global_optima.py,
%%  balance_lemma.py.
%%
%%  STATUS (2026-08-06).  REFUTED: the leximax objective (21 instances whose
%%  optima are ALL bad) and, again, the utilitarian one (50).  REFUTED: every
%%  lexicographic tie-break chain tried -- none makes all optima good, the best
%%  (balance then leximax) failing on 4 of 1,082 instances.  OPEN, with strong
%%  evidence: the BALANCE LEMMA, that every instance has a good BALANCED
%%  allocation -- zero failures over 1,082 instances and 466,560 balanced
%%  partitions, EF-free instances included.  Its value is that it restricts
%%  Conjecture conj:main to a far more rigid class of witnesses without any
%%  claim about local minima, which is what blocked Approach 7.
%%  NOTE for the algorithmic clause: minimising sum |B_i|^2 is trivial, but the
%%  lemma does not say WHICH balanced partition is good, so it is not by itself
%%  an algorithm.  See the last subsection.
%% ==========================================================================

\section{Approach 8: the global route}
\label{sec:approach8}

\subsection{Why global}

Approach 7 reduced Conjecture~\ref{conj:main} to a descent lemma, but
Remark~\ref{rem:descent-stronger} records that the descent lemma is strictly
stronger than the conjecture: it forbids bad \emph{local} minima of $\Psi$,
whereas Conjecture~\ref{conj:main} only concerns the \emph{global} minimum. The
global route drops the extra strength. It asks:

\begin{quote}
  is there an objective whose global optimum is always a good allocation?
\end{quote}

Such a statement would give existence with no reference to local search, and the
polynomial-time clause becomes a separate question about computing the optimum.
All objectives below are evaluated over partitions in canonical form
(Definition~\ref{def:canonical}), so every candidate is envy-freeable and
$\pathw{\alloc}$ is finite.

\subsection{What is refuted}

For an objective $f$, two questions differ: are \emph{all} $f$-optima good (a
usable theorem, since then any optimum may be returned), or is merely
\emph{some} $f$-optimum good (which needs a tie-break before it is a theorem)?
Over $1{,}760$ instances at $n \le 5$, $m \le 7$:

\begin{center}
\begin{tabular}{llrr}
\toprule
objective & & all optima good & some optimum good \\
\midrule
$\sum_i \cost_i(B_i)$ & utilitarian    & $389$ fail & $50$ fail \\
sorted cost vector    & leximax        & $54$ fail  & $21$ fail \\
\midrule
$\max_i \cost_i(B_i)$ & egalitarian    & $382$ fail & $0$ \\
$\sum_i \abs{B_i}^2$  & balance        & $68$ fail  & $0$ \\
$\max_i \cost_i - \min_i \cost_i$ & spread & $201$ fail & $0$ \\
\bottomrule
\end{tabular}
\end{center}

The utilitarian row restates the refutation of Section~\ref{sec:approach6}. The
leximax row is new and worth recording, because leximax is the objective the
goods literature reaches for first: on $21$ instances \emph{every} leximax
optimum is bad, so no tie-break can rescue it. Egalitarian, balance and spread
all survive in the weaker sense.

A tie-break chain would upgrade a ``some'' to a theorem, and twelve were tried.
None succeeds:

\begin{center}
\begin{tabular}{lr@{\qquad}lr}
\toprule
chain & failures & chain & failures \\
\midrule
bal.                       & $75$ & bal.\ $\to$ lex.          & $4$ \\
bal.\ $\to$ util.          & $21$ & bal.\ $\to$ lex.\ $\to$ util. & $4$ \\
bal.\ $\to$ egal.          & $30$ & bal.\ $\to$ egal.\ $\to$ lex.\ $\to$ util. & $4$ \\
bal.\ $\to$ egal.\ $\to$ util. & $8$ & egal.\ $\to$ bal.    & $48$ \\
bal.\ $\to$ util.\ $\to$ egal. & $8$ & egal.\ $\to$ bal.\ $\to$ util. & $9$ \\
bal.\ $\to$ spread         & $15$ & bal.\ $\to$ \#at max $\to$ util. & $29$ \\
\bottomrule
\end{tabular}
\end{center}

Adding components helps only down to $4$ failures out of $1{,}082$ and then
stops; the last three chains are progressively finer and all fail on the same
residue. So the good allocations are not picked out by any natural cost-based
ordering. Note also that appending ``$\#$ agents at maximum cost'' actively
\emph{breaks} the objective, turning $0$ ``some'' failures into $21$.

\subsection{What survives: the balance lemma}

The balance objective is special because its optima are describable outright:
$\sum_i \abs{B_i}^2$ is minimised exactly by the \emph{balanced} partitions,
those with all bundle sizes within $1$ of one another, i.e.\ every bundle of size
$\lfloor m/n \rfloor$ or $\lceil m/n \rceil$. So its ``some optimum is good'' row
reads as a structural restriction:

\begin{conjecture}[Balance lemma]
\label{conj:balance}
Every instance with negative dichotomous valuations admits a good allocation
whose partition is balanced.
\end{conjecture}

\begin{proposition}
\label{prop:balance-implies-main}
Conjecture~\ref{conj:balance} implies Conjecture~\ref{conj:main}, except for its
polynomial-time clause.
\end{proposition}

\begin{proof}
Immediate: a good allocation is one with $\max_i \pathw{\alloc}(i) \le 1$, which
by Observation~\ref{obs:integrality} and
Proposition~\ref{prop:zero-coordinate} yields $\optsubsidy \in \set{0,1}^n$ with
total at most $n-1$.
\end{proof}

\begin{remark}[Relation to Target G-bal, and a correction]
\label{rem:balance-vs-targetGbal}
Conjecture~\ref{conj:balance} was arrived at here by scanning global objectives,
but the restriction to cardinality-balanced partitions is not new to this
section: it is Definition~\ref{def:targetGbal}, introduced in
Section~\ref{sec:approach6} on the goods side of the size-shift involution. The
two are not the same statement, and the direction matters.
Proposition~\ref{prop:targetG-implies} bounds $\pathw{\alloc}(u)$ by
$q_u - \min_v q_v$ in one direction only, so Target G-bal implies
Conjecture~\ref{conj:balance} but not conversely --- and the gap is real, not
merely unproved. Over the balanced partitions enumerated in
\texttt{update\_22}, $666$ are good while carrying $q$-spread at least $2$,
whereas none has $q$-spread at most $1$ without being good. So
Conjecture~\ref{conj:balance} is a strict weakening of Target G-bal, and to that
extent a better target; but the idea of restricting to balanced partitions was
already in hand, and this section should not be read as introducing it.
\end{remark}

Conjecture~\ref{conj:balance} is stronger than Conjecture~\ref{conj:main}, and it
is worth being explicit about why that strengthening is of a more useful kind
than the one in Approach 7. The descent lemma strengthens the conjecture by
constraining the \emph{optimisation landscape} --- it must have no bad local
minimum --- which is a statement about every partition. The balance lemma
strengthens it by constraining the \emph{witness}, and a constrained witness is
extra structure a proof may use: on a balanced partition every bundle has one of
two sizes, $m$ is nearly determined by $n$ and the common size, and the pigeonhole
arguments that drive Example~\ref{ex:lowerbound} and the small-bundle theorem of
Section~\ref{sec:solved} apply directly.

\begin{center}
\begin{tabular}{lrr}
\toprule
family & instances & no good balanced allocation \\
\midrule
random, $n \le 5$, $m \le 7$                & $1{,}760$ & $0$ \\
random, $n \le 5$, $m \le 8$                & $1{,}082$ & $0$ \\
EF-free only                                & $30$      & $0$ \\
\midrule
balanced partitions examined                & \multicolumn{2}{r}{$466{,}560$} \\
\bottomrule
\end{tabular}
\end{center}

The two approaches also fit together rather than competing. The second
coordinate of $\Psi$ in Definition~\ref{def:potential} is exactly
$\sum_i \abs{B_i}^2$, so descent under $\Psi$ drives the partition towards
balance; Conjecture~\ref{conj:balance} asserts that the balanced region is where
the good allocations are. That the empirically successful potential and the
empirically successful global objective agree on the same quantity is weak
evidence that balance, not cost, is the operative structure.

\subsection{The induction the balance lemma invites, and why it does not close}
\label{sec:approach8-induction}

Restricting the witness suggests building it one chore at a time. Three
statements of increasing strength were tested; the first holds, and the two that
would prove anything fail.

\begin{proposition}[Incremental structure, empirical]
\label{prop:inc}
Write $(\mathrm{INC})$ for: there is a good balanced allocation $\alloc$, with its
assignment fixed, and an ordering $g_1,\dots,g_m$ of the chores such that
$\alloc$ restricted to $\set{g_1,\dots,g_k}$ is good for every $k$. Over
$1{,}125$ instances, $(\mathrm{INC})$ never failed --- not even when the
assignment was forbidden to change between steps.
\end{proposition}

Restrictions of dichotomous cost functions to a subset of chores are again
dichotomous, so each prefix is a legitimate smaller instance and
$(\mathrm{INC})$ has exactly the shape of an induction on $m$. It is not one.
$(\mathrm{INC})$ \emph{presupposes} a good balanced allocation and then asserts
it can be reached; it says nothing about existence, and the induction does not
close, because the good allocation supplied by the hypothesis on $m-1$ chores
need not be the restriction of any good allocation on $m$.

The statement that would prove existence has to be self-contained:

\begin{quote}
  $(\mathrm{NS})$ \; if $\alloc$ is a good allocation of a proper subset $T$ of
  the chores, some chore outside $T$ can be added to some bundle leaving it good.
\end{quote}

$(\mathrm{NS})$ implies Conjecture~\ref{conj:main} by induction from the empty
allocation, which is good since all costs are $0$, and gives an $O(m^2n^4)$
algorithm with no backtracking. It is false. Over $502{,}528$ good partial
allocations, $846$ admit no good one-chore extension, and $150$ resist even when
re-matching is permitted. The first witness is
\[
  \bundle{} = (\emptyset,\; \set{g_1,g_2,g_3},\; \emptyset),
  \qquad g_4 \text{ unassigned},
\]
three chores piled on one agent with two agents empty --- a state a balanced
construction would never reach. That suggests the repair: keep the partial
allocation balanced, inserting only into minimum-size bundles, which preserves
balance and would prove the balance lemma and Conjecture~\ref{conj:main}
together.

\begin{proposition}[The balanced induction fails, and fails worse]
\label{prop:nsb-fails}
Write $(\mathrm{NSB})$ for $(\mathrm{NS})$ restricted to balanced partial
allocations with insertion into a minimum-size bundle. Over $236{,}421$ good
balanced partial allocations, $11{,}723$ are stuck --- $5.0\%$, against $0.17\%$
for the unrestricted $(\mathrm{NS})$. Allowing insertion into any
balance-preserving bundle rather than a minimum-size one changes nothing: the
same $11{,}723$ states are stuck.
\end{proposition}

So balance does not rescue the induction; it makes it worse, by removing exactly
the freedom --- piling chores onto one agent for a step --- that lets greedy
insertion escape. Note that this does not contradict
Proposition~\ref{prop:inc}: a good insertion order exists, but it cannot be found
by extending an arbitrary good balanced partial allocation. The order must be
chosen with foresight.

\begin{remark}[The pattern across Approaches 7 and 8]
\label{rem:pattern}
Three independent constructive routes have now failed in the same way. The
descent lemma is true in every test but its witnessing transfer is not local to
the envy (Section~\ref{sec:approach7-obstruction}); $(\mathrm{INC})$ is true in
every test but the insertion order is not locally determined; and no
lexicographic objective picks out the good allocations. In each case the object
exists and no rule computable from local data finds it. That is weak evidence
that Conjecture~\ref{conj:main}, if true, is not provable by exhibiting a
construction, and that the search should turn to non-constructive existence
arguments --- counting, averaging, or a fixed-point or topological argument ---
with the algorithmic clause pursued separately through
Section~\ref{sec:approach6}, whose algorithm is polynomial by construction and
has never failed.
\end{remark}

\subsection{The non-local route: is there room, and can anything reach it?}
\label{sec:approach8-nonlocal}

Remark~\ref{rem:cancellation-four} closes the local and decompositional routes:
four independent attempts to bound $\pathw{\alloc}$ by a part of itself all came
in strictly worse than the truth. What remains is the non-local class --- counting,
averaging, probabilistic-method and parity arguments, which reason about all
allocations at once. Two questions decide whether that class is usable: whether
the good allocations are numerous enough to be reached by such an argument, and
whether any natural distribution or identity actually reaches them.

\begin{observation}[There is room]
\label{obs:counting-room}
Exhaustively over all $n^m$ allocations of $1{,}169$ instances with $n \le 5$,
$m \le 6$, the number of good allocations never fell below $11$, and on the $80$
EF-free instances --- the residual class, since Conjecture~\ref{conj:main} is
trivial where an exactly envy-free allocation exists --- never below $27$
(\texttt{update\_31/counting\_room.py}).
\end{observation}

So the good set is not a needle: an argument that merely shows it non-empty has
something to aim at, unlike the situation where only one witness exists and
nothing short of exhibiting it can work. But the two standard ways of reaching it
both fail.

\begin{proposition}[Parity fails]
\label{prop:parity-fails}
The number of good allocations has no fixed parity: over $737$ instances it was
even $469$ times and odd $268$ times. No Sperner- or Tucker-style identity forcing
an odd count can hold.
\end{proposition}

\begin{proposition}[No natural distribution concentrates]
\label{prop:no-distribution}
The fraction of allocations that are good is not bounded below. Under the uniform
distribution its minimum fell from $0.136$ at $n=3$, $m=4$ to $0.025$ at $n=m=5$.
Restricting to cardinality-balanced partitions does not repair this --- the
minimum fraction there was $0.042$ at $n=m=4$, \emph{worse} than uniform --- and
restricting to minimum-cost allocations gives a minimum fraction of $0$,
reconfirming Section~\ref{sec:approach6}'s refutation of that family.
\end{proposition}

The probabilistic method needs a distribution under which the good set has
probability bounded away from zero, and Proposition~\ref{prop:no-distribution}
rules out the three natural candidates; a counting proof needs an identity forcing
the count non-zero, and Proposition~\ref{prop:parity-fails} rules out the parity
identity. What is \emph{not} ruled out is a cleverer weighting: the fractions above
are for specific distributions, and the space of weightings is unbounded. But no
weighting suggested by the structure of the problem survives, and
Observation~\ref{obs:counting-room} shows the target is generous enough that if a
natural one existed it would likely have been among these.

\begin{remark}[Status of the technique classes]
\label{rem:technique-classes}
Taken with Remark~\ref{rem:cancellation-four}, every named technique class has now
been tried and refuted on this problem: constructive
(Sections~\ref{sec:approach7-obstruction}, \ref{sec:approach8-induction}),
decompositional (Remark~\ref{rem:cancellation}, four instances), transplanted from
the additive literature (Sections~\ref{sec:approach6-bdnsv-audit},
\ref{sec:approach6-r11-audit}), and now non-local. This is not a claim that
Conjecture~\ref{conj:main} is false --- it survives every test at every size --- nor
that it is unprovable. It is a claim about where a proof will not come from, and
it is the most useful thing this report currently establishes.
\end{remark}

\subsection{The algorithmic clause}
\label{sec:approach8-algorithmic}

Conjecture~\ref{conj:main} also asks for a polynomial-time algorithm, and
Proposition~\ref{prop:balance-implies-main} deliberately excludes it. The
position is as follows.

\begin{itemize}
  \item Minimising $\sum_i \abs{B_i}^2$ is trivial --- every balanced partition
    is optimal --- but that is exactly why the balance lemma gives no algorithm
    on its own: it does not say \emph{which} balanced partition is good, and
    there are exponentially many.
  \item The objectives that would name one are not obviously tractable.
    Minimising $\max_i \cost_i(B_i)$ is a makespan-type problem and NP-hard
    already for identical additive costs, so ``compute the egalitarian optimum''
    is not a polynomial algorithm even where it is correct.
  \item Two polynomial candidates remain and are unaffected by this section: the
    algorithm of Section~\ref{sec:approach6}, whose completeness is
    Conjecture~\ref{conj:algorithm-succeeds} and which has never failed in
    $389{,}215$ instances; and descent on $\Psi$, which
    Theorem~\ref{thm:descent-suffices} shows runs in $O(m^4 n^4)$ and is correct
    given Conjecture~\ref{conj:descent}.
\end{itemize}

So the division of labour is now explicit: the global route is the cleanest line
on \emph{existence}, and the algorithmic clause rests on Approach 6 or Approach
7. A proof of Conjecture~\ref{conj:balance} would settle the mathematical
content of Conjecture~\ref{conj:main} and leave a well-posed, separately
attackable computational question --- which is a better position than the
current one, where existence and computation are entangled in a single
conjecture that has resisted both.
